% ===========================================================================
%  91 — Decisões em aberto para o produto
% ===========================================================================
\chapter{Decisões em aberto para o produto}
\label{cap:decisoes}

O levantamento encontrou perguntas que o código não responde e que não são de
engenharia: são escolhas de produto. Cada uma tem o contexto, as opções e,
quando há, uma recomendação. Nenhuma foi decidida neste documento.

\newcommand{\decisaoaberta}[4]{%
  \begin{historia}{#1}{#2}{}
  \textbf{Contexto.} #3\par\smallskip
  #4
  \end{historia}}

\decisaoaberta{D01}{Quem cria o ano letivo?}{A tela \enquote{Ano letivo} diz
\enquote{O calendário é do contrato, não da escola}, mas a Coordenação da
escola consegue criar período (\id{H02.4.1}).}{\textbf{Opções:} (a) só o
Administrador do contrato cria, e a escola consulta; (b) a escola cria, e a
frase da tela muda. \textbf{Recomendação:} (a), porque um calendário por
escola quebra o relatório da rede.}

\decisaoaberta{D02}{A Coordenação da escola pode transferir estudante?}{O botão
\enquote{Transferir de escola} aparece só para o Administrador do contrato; a
Coordenação consegue o mesmo efeito importando o CPF (\id{H03.6.1}).}{\textbf{Opções:}
(a) mostrar o botão também para as coordenações; (b) manter só no contrato e
tratar a importação como caminho oficial da escola.}

\decisaoaberta{D03}{O texto livre da criança é só do professor em todas as
missões?}{No Rastro do Boato e no Manifesto da IA a gestão ainda vê o texto;
a versão nova de Quem mandou isso? já segue a regra --- a gestão vê que a regra
foi enviada, e só quem dá aula na turma lê o texto (\id{H07.1.4},
\id{H06.11.5}).}{\textbf{Recomendação:} alinhar todas à regra de \id{H07.1.4}
--- a gestão vê só \enquote{entregou / não entregou} ---, e tornar a regra
obrigatória para toda missão nova (\id{H06.10.1}).}

\decisaoaberta{D04}{A que horas o limite de IA recomeça?}{O contador de cada
pessoa zera à meia-noite no horário universal, que no Brasil é 21h
(\id{H09.4.1}). Um professor que planeja à noite pode ver o limite
\enquote{renascer} no meio do trabalho.}{\textbf{Opções:} (a) meia-noite no
fuso da rede; (b) manter e explicar na tela.}

\decisaoaberta{D05}{Como o estudante recupera a senha?}{Hoje só a coordenação
reseta. Estudante dos anos iniciais quase nunca tem e-mail
(\id{H01.4.2}).}{\textbf{Opções:} (a) recuperação por e-mail só para adultos, e
o estudante segue pela coordenação; (b) o professor também pode resetar a
senha dos estudantes da turma dele. Com o login único da rede (\id{H01.5.1}),
a pergunta pode deixar de existir.}

\decisaoaberta{D06}{A conta terá perfil e foto?}{O menu do avatar tem só
\enquote{Sair}, e isso foi deliberado.}{Confirmar que continua fora do escopo
--- foto de criança é dado sensível e não traz valor pedagógico evidente.}

\decisaoaberta{D07}{O estudante terá um histórico fora das missões?}{Hoje ele
vê só a grade da turma; não há portfólio nem \enquote{o que já fiz}.}{Decidir
se entra no produto e, se entrar, o que ele mostra sem virar nota.}

\decisaoaberta{D08}{Toda missão terá relatório individual?}{Só unooGPT, Vila
do Lume e Instanooria têm (\id{H07.1.3}).}{Decidir se passa a ser exigência
da esteira de missões novas (\id{H06.10.1}) ou se fica a critério de cada
missão.}

\decisaoaberta{D09}{Por quanto tempo guardamos o dado de criança?}{Não há prazo
definido, e a Política de Privacidade declara isso como pendência
(\id{H12.3.1}).}{Definir, com o jurídico, um prazo por tipo de dado --- é
pré-requisito para o desenvolvimento da rotina de eliminação.}

\decisaoaberta{D10}{Qual é o modelo das próximas missões?}{O Instanooria tem
\enquote{atividade} configurável (tema, geração, revisão, etapas), e a versão
nova de Quem mandou isso? tem o pacote de mensagens que o professor prepara e
aprova para a turma, que é a atividade dela (\id{H06.11.1}); as outras seguem
\enquote{liberar e conduzir}. Em 25 de setembro de
2026, com a evidência pedagógica (\id{E14}), ficou decidido que cada missão
define as próprias atividades e que o unooGPT e o UnooTube serão refeitos com
atividades numa próxima onda; o modelo das demais continua em aberto.}{Decidir
o modelo padrão: a atividade configurável dá mais controle ao professor e custa
bem mais para construir e para preparar a aula --- e é por atividade que a
evidência pedagógica é registrada.}

\decisaoaberta{D11}{A rede verá algum agregado do bem-estar?}{Hoje não existe
nenhuma visão de gestão, por desenho (\id{H08.3.4}).}{\textbf{Opções:} (a)
nunca; (b) um agregado anônimo, com grupos mínimos e termo específico com a
rede. Qualquer caminho diferente de (a) precisa do jurídico antes.}

\decisaoaberta{D12}{A liberação com data entra antes de novembro?}{Hoje o
professor libera na hora da aula, e é fácil esquecer (\id{H04.2.3}).}{Decidir
a prioridade.}

\decisaoaberta{D13}{Gravação de sessão por terceiros em tela de criança}{Uma
ferramenta externa de gravação de sessão é usada nas telas
(\id{H12.2.1}).}{\textbf{Opções:} (a) declarar no termo de tratamento de dados
de cada rede; (b) retirar das telas de estudante. \textbf{Recomendação:} (b)
nas telas de menores.}

\decisaoaberta{D14}{No pacote pronto de Quem mandou isso?, imagem e link
equilibrados?}{No material de origem do pacote pronto \enquote{O Mistério do
Passeio ao Museu}, os dois links eram falsos --- e a turma aprenderia que link
é sempre mentira, o atalho que a missão existe para desmontar. Até a resposta,
vale o equilíbrio: a plataforma exige que a imagem e o link tenham uma
verdadeira e uma falsa cada, e no pacote pronto a quinta mensagem (verdadeira)
passou a ser link e a sétima (falsa) passou a ser texto, sem mudar nenhum texto
(\id{H06.11.1}).}{\textbf{Opções:} (a) manter o equilíbrio, para o pacote pronto
e para o gerado; (b) voltar ao material original e aceitar que o formato
entregue a resposta. \textbf{Recomendação:} (a).}
